\clearpage\thispagestyle{empty}
\pdfbookmark[0]{Quellen zur Besprechung von Bodes Astronomischem Jahrbuch für 1822}{bode1822-sources}
\begin{center}\Large\bfseries Quellen zur Besprechung von Bodes Astronomischem Jahrbuch für 1822\end{center}

Die Übertragung folgt Carl Friedrich Gauß’ Besprechung in den \emph{Werken}, Bd.~VI (Göttingen, 1874), S.~619--622. Der Erstabdruck erschien in den \emph{Göttingischen gelehrten Anzeigen}, Stück~23, 7.~Februar 1820, S.~225--232. Johann Elert Bodes Jahrbuch für 1822 erschien 1819 in Berlin beim Verfasser und in Kommission bei Ferdinand Dümmler.

Die in der Besprechung zitierten Druckfehler stehen auf \href{https://api.digitale-sammlungen.de/iiif/image/v2/bsb10538351_00245/full/1650,/0/default.jpg}{S.~239 des Jahrbuchs}.

KI-Übertragung aus dem historischen Deutsch in modernes Deutsch. KI-erstellter Satz.

\href{https://gdz.sub.uni-goettingen.de/id/PPN236007467}{Werkausgabe}. Erstdruck, S.~\href{https://images.p.sub.uni-goettingen.de/iiif/image/gdz:PPN319721507_1820:00000219/full/1650,/0/default.jpg}{225}, \href{https://images.p.sub.uni-goettingen.de/iiif/image/gdz:PPN319721507_1820:00000220/full/1650,/0/default.jpg}{226}, \href{https://images.p.sub.uni-goettingen.de/iiif/image/gdz:PPN319721507_1820:00000221/full/1650,/0/default.jpg}{227}, \href{https://images.p.sub.uni-goettingen.de/iiif/image/gdz:PPN319721507_1820:00000222/full/1650,/0/default.jpg}{228}, \href{https://images.p.sub.uni-goettingen.de/iiif/image/gdz:PPN319721507_1820:00000223/full/1650,/0/default.jpg}{229}, \href{https://images.p.sub.uni-goettingen.de/iiif/image/gdz:PPN319721507_1820:00000224/full/1650,/0/default.jpg}{230}, \href{https://images.p.sub.uni-goettingen.de/iiif/image/gdz:PPN319721507_1820:00000225/full/1650,/0/default.jpg}{231}, \href{https://images.p.sub.uni-goettingen.de/iiif/image/gdz:PPN319721507_1820:00000226/full/1650,/0/default.jpg}{232}. \href{https://api.digitale-sammlungen.de/iiif/image/v2/bsb10538351_00003/full/1650,/0/default.jpg}{Bodes Jahrbuch, Bayerische Staatsbibliothek}.

\clearpage
